\documentclass[level=1,teacher]{kidscode}
%% small helpers for this guide
\newcommand{\wk}[1]{\textbf{W#1}}
\newenvironment{guidetable}[1]{\par\noindent\small\renewcommand{\arraystretch}{1.3}%
  \longtable{#1}}{\endlongtable\par}
%% \fitwidth{wide content}: shrink only if it is wider than the line (dyslexia edition)
\newsavebox{\fitbox}
\newcommand{\fitwidth}[1]{\sbox{\fitbox}{#1}\ifdim\wd\fitbox>\linewidth\resizebox{\linewidth}{!}{\usebox{\fitbox}}\else\usebox{\fitbox}\fi}
\begin{document}
\kctitle{Teacher Guide}{Level 1: Introductory. Everything you need before, during and between sessions}
\tableofcontents

%% =====================================================================
\section{Welcome}
Thank you for teaching this course. You do not need to be a programmer or a trained teacher to run it
well. Every session is written out for you: what to print, what to open, what to say, what to do when
things go wrong, and what to cut if time runs short. If you can read a lesson plan for ten minutes the
day before, you are ready.

Level 1 is for complete beginners aged about 9--11 (US Grades 4--5). In weeks 1--6 children learn the big ideas of
programming by snapping blocks together in \textbf{Scratch}; in weeks 7--12 they type the same ideas as
real \textbf{Python} code, drawing pictures with the \enquote{turtle}. The group's love of games and
visual art runs through everything: animation, pen art, a catch game, random art and a final
\enquote{Interactive Art Show}.

\begin{teachernote}[Three things that matter most]
\begin{enumerate}[itemsep=2pt]
  \item \textbf{Bugs are normal.} Celebrate them. Say \enquote{Great, a bug! Let's catch it} and mean it.
        Your calm when things break is the most important thing you teach.
  \item \textbf{Hands off the keyboard.} Point, ask and suggest; let the child type and click.
        (\enquote{What did you want it to do? What did it do instead?})
  \item \textbf{Creativity, working code and explaining are equally valued.} A simple program a child
        can explain beats a complicated one they copied.
\end{enumerate}
\end{teachernote}

\begin{teachernote}[The math this level uses]
Everything fits Grade 4--5 math: turns use degrees and the rule \enquote{turn $= 360 \div$ number of sides}
(Grade 4), and places use $x$ and $y$ like a map grid (Grade 5). One idea comes early: the middle of the
Scratch stage and of the turtle window is $(0, 0)$, so places to the left or below use minus numbers such as
$-100$, which school math reaches in Grade 6. Children only need \enquote{a minus sign means the other way}:
the workbook says so in week 1 (Scratch) and week 8 (Python). In Python, \texttt{*} means $\times$ and
\texttt{/} means $\div$ (workbook chapter 8), and lists count from 0 (chapter 9).
\end{teachernote}

\subsection{How the curriculum works}
There are three levels of 12 weekly 60-minute sessions, plus two floating \enquote{buffer} sessions per
level. Inside each level the rhythm is the same: weeks 1--5 new ideas, \textbf{week 6 review and
checkpoint}, weeks 7--9 new ideas, weeks 10--11 final project, \textbf{week 12 showcase and exit
assessment}. Big ideas come back again and again (the \enquote{spiral}, see the last section): each
concept is first \textbf{I}ntroduced, later \textbf{R}evisited and then \textbf{E}xtended.

\begin{guidetable}{@{}>{\bfseries\RaggedRight}p{4.1cm}>{\RaggedRight}p{2.2cm}>{\RaggedRight\arraybackslash}p{\dimexpr\linewidth-6.3cm-4\tabcolsep}@{}}
\toprule
Document & Who & What it is for\\
\midrule\endhead
\bottomrule\endfoot
Teacher Guide (this) & You & Setup, how sessions run, assessment, troubleshooting, pacing.\\
Lesson Plans & You & One plan per session (and B1, B2): minute-by-minute, with scripts to read aloud.\\
Student Workbook & Students & One chapter per week: explanation, pictures, in-class exercises, new words, reflection.\\
Homework & Students & 15--20 minutes a week: Core, optional Challenge and Creative, self-check.\\
Answer Key & You only & Answers to workbook, homework and assessments, with scoring notes.\\
Assessments & Students & Week-1 survey, week-6 checkpoint, week-12 exit assessment.\\
Project Brief / Project Teacher Guide & Students / you & The final project, plan sheet and rubric / how to run it, reference solution, showcase.\\
Extras & Everyone & Cheat sheet, glossary, progress tracker, notes for families, certificate.\\
\end{guidetable}

\textbf{The code folder.} Each in-class exercise lives in \texttt{grade-4-6/1-introductory/\allowbreak{}exercises/\allowbreak{}weekNN/\allowbreak{}ex-name/\allowbreak{}},
homework coding in \texttt{hw-name/}. A Python exercise folder has a \texttt{README.md} (the task),
\texttt{starter.py} (what students open and edit), \texttt{check\_my\_work.py} (students press Run on it
to get friendly automatic feedback) and a commented \texttt{solution.py} for you. Scratch and Terminal exercises have a
\texttt{README.md} (what to build, with a checklist) and a \texttt{solution.md} for you, which lists every script block by block. Every PDF also comes in a
\textbf{dyslexia-friendly edition} (clearer font, larger spacing, cream paper color).

\subsection{Level 1 goals}
By the end of Level 1 the student can\dots{} (CSTA = CSTA K--12 Computer Science Standards, 2017; level 1B is grades 3--5.)

\begin{guidetable}{@{}>{\bfseries}p{1.3cm}>{\RaggedRight}p{\dimexpr\linewidth-5.6cm-4\tabcolsep}>{\RaggedRight\arraybackslash}p{4.3cm}@{}}
\toprule
Code & Goal & CSTA\\
\midrule\endhead
\bottomrule\endfoot
L1.G1 & write precise step-by-step instructions (algorithms) and put them in the right order & 1B-AP-08, 1B-AP-10\\
L1.G2 & use loops to repeat actions and make patterns & 1B-AP-10\\
L1.G3 & use events and conditionals to make programs react & 1B-AP-10\\
L1.G4 & use variables (score, size, answers) and random numbers & 1B-AP-09\\
L1.G5 & break a problem into parts and make reusable blocks/functions & 1B-AP-11, 1B-AP-12\\
L1.G6 & find and fix bugs using a strategy & 1B-AP-15, 1B-CS-03\\
L1.G7 & work in a pair (driver/navigator) and explain their code & 1B-AP-16, 1B-AP-17, 1B-IC-20\\
L1.G8 & stay safe and fair online: privacy, credit, kindness, accessible design & 1B-NI-05, 1B-AP-14, 1B-IC-21, 1B-IC-19\\
L1.G9 & read and write simple text-based Python programs in Thonny & 1B-AP-10, 1B-AP-12\\
\end{guidetable}

\subsection{The 12-week roadmap}
\begin{guidetable}{@{}>{\bfseries}p{0.7cm}>{\RaggedRight}p{3.5cm}>{\RaggedRight}p{\dimexpr\linewidth-10.4cm-8\tabcolsep}>{\RaggedRight}p{1.5cm}>{\RaggedRight\arraybackslash}p{4.7cm}@{}}
\toprule
Wk & Title & Big idea & Tool & Highlights\\
\midrule\endhead
\bottomrule\endfoot
1 & Robots Need Instructions & A program is a list of exact steps in order. & Scratch & Parent-robot drawing; class agreements; nicknames.\\
2 & Loop-de-Loop & Loops repeat steps so we don't have to. & Scratch & Dance loops; pen shapes ($360 \div n$); animation.\\
3 & When This Happens\dots & Events start code; conditionals choose. & Scratch & Simon-says-IF; \enquote{two stars and a wish}.\\
4 & Keeping Score & Variables are labeled boxes that remember. & Scratch & Catch game with score and random; credit and remix.\\
5 & Messages and My Blocks & Break big programs into named parts. & Scratch & Recipe cards; broadcast; My Blocks; accessible colors.\\
6 & Review Studio and Checkpoint & Show what you know; get ready for text code. & Scratch & Bug-hunt relay; Scratch-to-Python card sort; \textbf{checkpoint (15 min)}.\\
7 & Hello, Python! & Text code is blocks written as words, and Python is picky. & Thonny & First turtle square; typo hunt; saving files safely.\\
8 & Loops in Python & \texttt{for} loops + indentation = repeat. & Thonny & Indent the line-up; polygons; nested loops (light).\\
9 & Variables, Input and Choices & Programs ask, remember and decide. & Thonny & Choose-your-path floor maze; \texttt{input}/\texttt{int}; questions a program should not ask.\\
10 & Random Art and My Own Commands & \texttt{random} surprises; \texttt{def} makes commands. & Thonny & Function machine; \textbf{project kickoff (M1 plan)}.\\
11 & Project Studio & Build, test and improve in small steps. & Thonny & Paper prototype; partner-swap bug hunt; \textbf{M2, M3}.\\
12 & Showcase and Celebration & Share, explain, reflect. & Thonny & Family gallery; \textbf{exit assessment (20 min)}; certificates.\\
B1 & Unplugged Games Day & (any time) & Paper & Pixel-art grids, algorithm relay race, debugging detective cards.\\
B2 & Free-Build and Remix Studio & (any time) & Either & Remix a past exercise; gallery share.\\
\end{guidetable}

%% =====================================================================
\section{Before you start: setup checklist}
Allow about 90 minutes, once, a week or more before session 1. You need the internet only for these
downloads; after that everything works offline.

\subsection{1. Install Scratch (offline app) on every Mac}
\begin{checklist}
  \item Open Safari and go to \texttt{scratch.mit.edu/\allowbreak{}download}.
  \item Under \enquote{Install the Scratch app}, choose \textbf{macOS}. Either click the Mac App Store
        link (uses your own Apple ID) or choose the direct download and open the \texttt{.dmg} file.
  \item If you used the \texttt{.dmg}: drag \textbf{Scratch 3} into the \textbf{Applications} folder.
  \item Open Scratch once. If macOS says it cannot check the app, Control-click the app, choose
        \menu{Open}, then \menu{Open} again.
  \item Test it offline: turn Wi-Fi off, open Scratch, drag a \enquote{move 10 steps} block and click it.
        No account and no sign-in are needed, and the app has no \enquote{Share} button.
  \item Saving: \menu{File > Save to your computer}. Put projects in a folder called
        \texttt{My Scratch} in Documents. Opening: \menu{File > Load from your computer}.
\end{checklist}

\subsection{2. Install Thonny on every Mac}
\begin{checklist}
  \item Go to \texttt{thonny.org} and download the macOS installer (\texttt{.pkg}). Thonny includes its
        own Python, so nothing else is needed.
  \item Open the \texttt{.pkg} and click through the installer.
  \item Open Thonny. On first start choose your language and \textbf{Standard} settings.
  \item Test: type \texttt{print("hello")} in the editor, press the green \menu{Run} button (or \keys{F5}),
        save as \texttt{test.py} when asked, and see \texttt{hello} in the Shell at the bottom.
\end{checklist}

\subsection{3. Copy the course folder to every Mac}
\begin{checklist}
  \item On your own computer, make the \textbf{student copy} of the course: in the course folder run
        \texttt{python3 tools/student\_bundle.py -{}-level 1} (or \texttt{make student-bundle}). It creates
        \texttt{build/\allowbreak student-bundle/\allowbreak kids-coding/}: every exercise folder \emph{without} the solutions, the
        project starter files, the small \texttt{testing} helper that \texttt{check\_my\_work.py} needs, and a
        \texttt{START-HERE.txt} for the child.
  \item Copy that \texttt{kids-coding} folder (USB stick, AirDrop or a shared drive) into each child's
        \textbf{Documents} folder. Do not copy only an exercise folder: \texttt{check\_my\_work.py}
        needs the \texttt{testing} folder further up. Do not give children the whole course folder either:
        it contains every answer key, reference solution and teacher guide. Keep that on \textbf{your} Mac.
  \item Do not put the printed answer key, lesson plans or this guide where students work.
\end{checklist}

\subsection{4. Install the checking helper (pytest) in Thonny}
\begin{checklist}
  \item In Thonny: \menu{Tools > Manage packages\dots}
  \item Type \texttt{pytest} in the box, click \menu{Search on PyPI}, click \textbf{pytest} in the results,
        then \menu{Install}. Wait for it to finish and close the window.
  \item Test: open \texttt{grade-4-6/1-introductory/\allowbreak{}exercises/\allowbreak{}week07/\allowbreak{}ex-first-square/\allowbreak{}check\_my\_work.py} and press
        \menu{Run}. You should see a box of \texttt{[PASS]} / \texttt{[TRY AGAIN]} lines (the untouched
        starter has only one side, so some checks say TRY AGAIN: that proves the checker works).
\end{checklist}

\subsection{5. Print}
\begin{guidetable}{@{}>{\bfseries\RaggedRight}p{3.3cm}>{\RaggedRight\arraybackslash}p{\dimexpr\linewidth-3.3cm-2\tabcolsep}@{}}
\toprule
When & What to print\\
\midrule\endhead
\bottomrule\endfoot
Before week 1 & One per student: Student Workbook (or print each week's chapter the week before), Cheat Sheet (laminate it if you can), Glossary, Progress Tracker, the week-1 survey from Assessments. One per family: Notes for Families. For you: this guide and the Lesson Plans (keep the Answer Key in your folder).\\
Every week & The lesson plan for that week (if you did not print them all), the homework set, any cards or sheets listed under \enquote{Materials} in the plan.\\
Week 6 & The mid-level checkpoint (Assessments).\\
Week 10 & The Project Brief (one per student). The Project Teacher Guide for you.\\
Week 12 & The exit assessment, feedback cards (last page of the Project Teacher Guide), one certificate per student (write names by hand).\\
\end{guidetable}
Choose the standard or the dyslexia-friendly edition per student. Both have the same page content, so
students can use different editions side by side.

\subsection{6. Room setup}
\begin{itemize}[itemsep=2pt]
  \item One Mac per \textbf{pair}, with two chairs side by side and a real mouse (dragging Scratch blocks
        with a trackpad is hard for small hands). Four students need two Macs; a third Mac for your
        demonstrations is useful but not required (students can gather around one).
  \item A clear table for unplugged activities, away from the computers. Drinks and snacks away from keyboards.
  \item A visible timer (a kitchen timer or a phone) for driver/navigator swaps.
  \item Pencils, paper, sticky notes, scissors, and a box for printed cards that you will reuse.
  \item A spot on the wall for the class agreements you write together in week 1.
\end{itemize}

\subsection{7. Accessibility settings}
\begin{itemize}[itemsep=3pt]
  \item \textbf{Dyslexia-friendly PDFs:} every document has a second edition (in the \texttt{dyslexia}
        PDF folder) with the Lexend font, larger text, extra spacing and a cream background. Offer it to any
        student who finds reading tiring; no diagnosis needed.
  \item \textbf{Bigger code in Thonny:} \menu{View > Increase font size} (or \keys{Cmd}+\keys{+}); repeat
        until comfortable. \menu{View > Decrease font size} goes back.
  \item \textbf{Bigger blocks in Scratch:} use the round \enquote{+} zoom button at the bottom right of the
        code area; the stage has a full-screen button above it.
  \item \textbf{Color filters on the Mac:} \menu{System Settings > Accessibility > Display > Color Filters}
        can help students who see colors differently. Also there: a larger pointer
        (\menu{Pointer > Pointer size}) and \menu{Zoom}. All course materials already use the color-blind-safe
        Okabe--Ito colors and never rely on color alone.
  \item \textbf{Reading support:} read instructions aloud for the whole group as a habit; accept spoken
        answers in assessments.
\end{itemize}

%% =====================================================================
\section{The 60-minute session}
Every session, including the buffers, has the same five parts. Children quickly learn the rhythm, which
makes the hour calmer and easier to run.

\begin{timeplan}
  \timeslot{0--8}{Warm-up}{An unplugged game, a quick quiz or \enquote{show me your homework}. Gets bodies and brains switched on and connects to last week.}
  \timeslot{8--18}{Concept}{The one new idea of the day, in kid language, with an analogy. Read the teacher script aloud; ask questions back.}
  \timeslot{18--28}{Guided demo}{You build the example step by step while students watch and predict (\enquote{What will happen when I click?}). Make one deliberate mistake and fix it together.}
  \timeslot{28--52}{Hands-on build}{Pairs work through the workbook exercises (driver/navigator, swap every 8 minutes), with a bug hunt. The last 5--6 minutes are creative free-build.}
  \timeslot{52--60}{Wrap-up}{Show-and-tell (each pair shows one thing), reflection question, exit ticket, progress tracker, hand out homework.}
\end{timeplan}

\begin{tip}[Making the hour work]
\begin{itemize}[itemsep=1pt]
  \item Start the computers and open the files \textbf{before} the children arrive (the materials box in every plan says what).
  \item Protect the wrap-up. When the build runs over, cut free-build first, never the wrap-up.
  \item A 1-minute stretch break around minute 30 (from week 2) helps everyone. Use the 20-20-20 eye rule.
  \item Every plan ends with \enquote{If you are short on time}: what to cut without losing the main idea.
\end{itemize}
\end{tip}

%% =====================================================================
\section{Pair programming}
Students work in pairs at one computer. This is how professional programmers often work, and it
doubles the thinking while halving the frustration.

\begin{pairroles}[The two roles]
\textbf{Driver:} has the mouse and keyboard; types and clicks; says out loud what they are doing.\\
\textbf{Navigator:} reads the instructions, thinks one step ahead, spots bugs, suggests what to try.
Never grabs the mouse: points at the screen and uses words.\\
\textbf{Swap} every 8 minutes when the timer rings, even in the middle of something. Saying
\enquote{Swap!} out loud and changing seats makes it a game.
\end{pairroles}

\begin{guidetable}{@{}>{\bfseries}p{2.4cm}>{\RaggedRight\arraybackslash}p{\dimexpr\linewidth-2.4cm-2\tabcolsep}@{}}
\toprule
Students & How to run it\\
\midrule\endhead
\bottomrule\endfoot
1 student & You are the navigator. Resist taking the mouse. Once a session, swap: the student navigates while you drive (and make a small deliberate mistake for them to catch).\\
2 students & One pair. Seat them so both can see the screen equally.\\
3 students & A rotating trio: \textbf{driver}, \textbf{navigator} and \textbf{tester} (runs \texttt{check\_my\_work.py}, reads hints aloud, keeps the bug list). Rotate every 6--7 minutes.\\
4 students & Two pairs on two Macs. Change partners every few weeks so the two friends also work with others. Pair by temperament (calm with energetic) rather than \enquote{strong with weak}.\\
\end{guidetable}

\begin{tip}[Phrases worth teaching]
Navigator: \enquote{I think the next step is\dots}, \enquote{Can we run it to check?}, \enquote{Look at line 5.}
Driver: \enquote{I'm going to\dots}, \enquote{What do you think?}
Both: \enquote{Let's try your idea first, then mine.}
If one child dominates, give the quieter one the navigator role and the instruction sheet.
\end{tip}

%% =====================================================================
\section{The four threads}
Four kinds of activity appear in almost every session. Each has its own labeled box in the lesson plans.

\subsection{Bug hunts}
\begin{bughunt}[How a bug hunt works]
Students get a program that \emph{almost} works. They (1) \textbf{predict} what it should do, (2) \textbf{run}
it and describe what actually happens, (3) \textbf{find} the bug, (4) \textbf{fix} it, and (5) \textbf{say
the strategy} they used. Celebrate every catch (\enquote{Bug caught!}). Over the level we build a
strategy list on the wall: read the error message; check the last thing you changed; read the code aloud;
act it out; compare with a version that works; test one small part at a time.
\end{bughunt}

\subsection{Unplugged activities}
\begin{unplugged}[Why and how]
Unplugged means no computers: cards, movement, paper, acting. It makes abstract ideas physical
(a loop becomes a dance; indentation becomes standing inside a hoop), gives eyes a rest, and lets
students who are not yet confident with the mouse shine. Keep them short (5--10 minutes), give
everyone a role, and always end by linking back: \enquote{Where will we see this in code today?}
Buffer B1 is a whole unplugged games day.
\end{unplugged}

\subsection{Creative free-build}
\begin{freebuild}[Five minutes of \enquote{make it yours}]
The last 5--6 minutes of the hands-on build are open: change colors, add a sprite, make a pattern,
add a sound. Two rules: \textbf{save your working version first}, and \textbf{one idea at a time}.
Free-build is where the art and games interests show up, and it is where many students learn the most.
If a pair is behind, they still get at least 3 minutes of free-build: it keeps motivation high.
\end{freebuild}

\subsection{Reflection}
\begin{reflection}[Closing the loop]
Every wrap-up ends with one short reflection question (in the lesson plan and the workbook), for
example \enquote{What was the trickiest bug today?} or \enquote{Where else do you see loops in real life?}
Students answer aloud or write a line in the workbook, then check the progress tracker. Model it
yourself: \enquote{Today I learned that I should read the whole error message.}
\end{reflection}

%% =====================================================================
\section{Assessment plan}
Assessment in this course is \textbf{for learning}: it tells you what to teach next and lets children
see their own progress. There are no grades. We value three things equally: \textbf{creativity},
\textbf{working code} and the \textbf{ability to explain} it.

\begin{guidetable}{@{}>{\bfseries\RaggedRight}p{3.2cm}>{\RaggedRight}p{2.4cm}>{\RaggedRight\arraybackslash}p{\dimexpr\linewidth-5.6cm-4\tabcolsep}@{}}
\toprule
What & When & How to use it\\
\midrule\endhead
\bottomrule\endfoot
Interest and comfort survey & Week 1, 10 min & Not graded. Learn interests, worries and learning preferences; choose examples and first pairs. Keep it to compare with week 12.\\
Formative checks & Every session & Watch the build, use the exit-ticket question in the plan, and look at \texttt{check\_my\_work.py} results. Note one line per student.\\
Progress tracker & End of every session, 2 min & Students check \enquote{Getting there / Got it / Expert} for that week's \enquote{I can\dots} lines; you add a sticker.\\
Mid-level checkpoint & Week 6, 15 min & Scratch reading and predicting, unplugged sequencing, finding bugs. 9 points. Use the results to re-teach before Python (or use B1).\\
Final project rubric & Weeks 10--12 & Five rows (creativity, working code, ingredients, explaining, process) scored 1--4. See the Project Teacher Guide.\\
Exit assessment & Week 12, 20 min & Predict Python output, fix bugs, a short build, explain code. 14 points. Combine with the project to decide next steps.\\
\end{guidetable}

\textbf{The four levels} used everywhere (tracker, rubrics, assessments):
\faSeedling\ \textbf{Just starting} (needs help or an example),
\faShoePrints\ \textbf{Getting there} (does it with some help),
\faCheckCircle\ \textbf{Got it} (does it alone),
\faRocket\ \textbf{Expert} (does it alone, explains it, uses it in a new way).
Answers, scoring bands and teacher rubrics for all three assessments are in the Answer Key.

\begin{tip}[Valuing the right things]
Praise strategies and explanations, not speed. When a student shows you a program, ask
\enquote{How does this part work?} before you say \enquote{Well done}. A student who can explain a short
program is further along than one who copied a long one.
\end{tip}

%% =====================================================================
\section{Differentiation}
Every lesson plan has a support path (needs more time) and an extension path (finished early). General moves:

\begin{differentiation}
\supportpath Give a partly finished program and ask them to change numbers rather than start from blank;
read instructions aloud; reduce the number of exercises (the first one is always the core); pair with a
patient partner as navigator; use the cheat sheet; enlarge fonts; allow the Scratch version of Python tasks.
\extensionpath Ask \enquote{Can you make it react to a key? Use a variable? Make it random?}; the
stretch challenge in each plan; ask them to write a bug for others to find; let them become the
\enquote{bug-hunt helper} who asks questions (but does not touch other people's keyboards).
\end{differentiation}

\begin{guidetable}{@{}>{\bfseries}p{0.7cm}>{\RaggedRight}p{\dimexpr(\linewidth-0.7cm-4\tabcolsep)/2}>{\RaggedRight\arraybackslash}p{\dimexpr(\linewidth-0.7cm-4\tabcolsep)/2}@{}}
\toprule
Wk & \faArrowDown\ Support path & \faArrowUp\ Extension path\\
\midrule\endhead
\bottomrule\endfoot
1 & Parent-robot with only four command cards; build \enquote{when flag clicked, move, say} only. & Add a second sprite with its own script; write a robot algorithm for a harder drawing.\\
2 & Start from a working square; change only the numbers to get a triangle. & Repeat a square while turning to make a flower; animate with costumes.\\
3 & One event and one \textbf{if}; use a ready-made sprite. & \textbf{if-else} with two keys; a simple chase game.\\
4 & Score only (no random yet); one falling object. & Add lives and \enquote{game over} with an \textbf{if}; speed up as the score grows.\\
5 & One broadcast between two sprites; a My Block without inputs. & A My Block with two inputs (size and sides) that draws any polygon.\\
6 & Bug-hunt relay with hint cards; checkpoint read aloud, spoken answers. & Become the bug-hunt helper; build a broken project for others to fix.\\
7 & Change numbers in the given program; navigator spells words letter by letter; bigger font. & Draw their initials or a house with \texttt{penup}/\texttt{goto}.\\
8 & One loop (a square), then change \texttt{range} and the angle. & Nested loops for a spirograph; use \texttt{i} to grow a spiral.\\
9 & One question and an \texttt{if}/\texttt{else} only. & Three choices with \texttt{elif}; check the answer with \texttt{isdigit()}.\\
10 & Use a given function and change its arguments. & A function with three or more parameters; a function that calls another one.\\
11 & Build the project in Scratch (same rubric); fewer ingredients first, then add. & Add key presses with \texttt{onkey} (project brief extension).\\
12 & Use the demo-script card; practice with a partner; oral exit assessment. & Tell the story of a bug they fixed; co-host the showcase.\\
\end{guidetable}

%% =====================================================================
\section{Digital citizenship thread}
Every week has a short safety, fairness or wellbeing moment (a \enquote{Digital citizenship \& safety}
box in the lesson plan). Keep it to 2--3 minutes and make it a conversation, not a lecture.

\begin{guidetable}{@{}>{\bfseries}p{0.7cm}>{\RaggedRight}p{4.2cm}>{\RaggedRight\arraybackslash}p{\dimexpr\linewidth-4.9cm-4\tabcolsep}@{}}
\toprule
Wk & Theme & The key message\\
\midrule\endhead
\bottomrule\endfoot
1 & Class agreements; privacy & We agree how we treat each other. We use the offline app, no accounts, and nicknames: never real names or addresses in projects.\\
2 & Screen-time balance & Stretch breaks and the 20-20-20 eye rule. Coding is better with rested eyes.\\
3 & Kind feedback & \enquote{Two stars and a wish}: two things you like, one idea to improve.\\
4 & Credit and remix & Use the built-in library; say whose art, sounds and ideas you used.\\
5 & Accessibility & Color-blind-safe colors plus shapes and words; readable text sizes.\\
6 & Fun and fair & What makes a project fun and fair for everyone who uses it?\\
7 & Saving files safely & Where your work lives; sensible file names; save often.\\
8 & Asking for help well & Say what you expected and what happened instead.\\
9 & Private information & Questions a program should never ask: real name, age, address, school, passwords.\\
10 & Being an original artist & Remixing is fine when you give credit; make it your own.\\
11 & Helpful feedback & Giving and receiving feedback kindly during playtesting.\\
12 & Showcase safety & Share with family only, on the screen; nothing uploaded; no personal information on screen.\\
\end{guidetable}

%% =====================================================================
\section{Troubleshooting}
\begin{guidetable}{@{}>{\RaggedRight}p{5.0cm}>{\RaggedRight\arraybackslash}p{\dimexpr\linewidth-5.0cm-2\tabcolsep}@{}}
\toprule
\textbf{Problem} & \textbf{What to do}\\
\midrule\endhead
\bottomrule\endfoot
\multicolumn{2}{@{}l}{\textbf{\color{kcaccent}Scratch}}\\ \midrule
Scratch will not open (\enquote{cannot be opened because Apple cannot check it}). & In Applications, Control-click \textbf{Scratch 3}, choose \menu{Open}, then \menu{Open}. Only needed once.\\
The pen blocks are missing. & Click the \menu{Add Extension} button (bottom left), choose \textbf{Pen}. Each new project needs it again.\\
A sprite has vanished off the stage. & Click the sprite in the sprite list, set x and y to 0 in the boxes above it, and check \enquote{Show} is on.\\
A project was not saved / cannot be found. & \menu{File > Save to your computer} and choose the \texttt{My Scratch} folder in Documents. Open with \menu{File > Load from your computer}.\\
Nothing happens when the green flag is clicked. & The script needs a \enquote{when green flag clicked} hat on top, and the blocks must be snapped together (drag the stack: does it move as one?).\\
\midrule
\multicolumn{2}{@{}l}{\textbf{\color{kcaccent}Thonny and Python}}\\ \midrule
The Run button is gray or nothing happens. & A program is still running. Press the red \menu{Stop} button (or \keys{Ctrl}+\keys{F2}), then Run again.\\
The turtle window will not close, or the beachball spins. & Click the turtle window and close it, or press \menu{Stop} in Thonny. \texttt{turtle.done()} keeps the window open on purpose until you close it.\\
The turtle window does not appear. & It may be hidden behind Thonny: press \keys{Cmd}+\keys{Tab} or click the Python rocket icon in the Dock.\\
\texttt{Terminator} error after closing the window. & Harmless. Press Run again.\\
\texttt{AttributeError: module 'turtle' has no attribute 'Turtle'} & The student saved their file as \texttt{turtle.py} (or \texttt{random.py}). Rename it (e.g. \texttt{my\_square.py}) and delete any \texttt{turtle.py} in that folder.\\
\texttt{ModuleNotFoundError:\allowbreak{} No module named 'pytest'} & Install pytest: \menu{Tools > Manage packages\dots}, search \texttt{pytest}, Install.\\
\texttt{ModuleNotFoundError:\allowbreak{} No module named 'kidtest'} & Only part of the student copy was copied. Copy the whole \texttt{kids-coding} folder (with \texttt{testing}) and open the exercise from there.\\
\texttt{IndentationError} (\enquote{expected an indented block} or \enquote{unexpected indent}). & Lines inside \texttt{for}, \texttt{if} and \texttt{def} need 4 spaces; other lines none. Select lines and press \keys{Tab} (in) or \keys{Shift}+\keys{Tab} (out). Check the line above ends with a colon.\\
\texttt{SyntaxError} & Look at the line Thonny highlights \emph{and the line above it}: a missing \texttt{)} or \texttt{"} or \texttt{:} is the usual cause. \menu{View > Assistant} explains errors in friendlier words.\\
\texttt{NameError: name 'Print' is not defined} & Spelling or capitals: Python words are lower case (\texttt{print}, \texttt{forward}).\\
\midrule
\multicolumn{2}{@{}l}{\textbf{\color{kcaccent}check\_my\_work.py}}\\ \midrule
\enquote{I need a helper called pytest}. & Install pytest (setup step 4).\\
\enquote{I could not find any checks to run}. & \texttt{check\_my\_work.py} was moved out of its exercise folder, or the folder was copied without its \texttt{test\_\dots.py} file. Use the original folder.\\
All checks say TRY AGAIN, but the drawing looks right. & The checker reads \texttt{starter.py}. Did the student \menu{Save as\dots} under a new name? Copy the code back into \texttt{starter.py} and save.\\
\enquote{asked for more input than the test gave it}. & The program asks more questions than the task said. Remove the extra \texttt{input()} or check the README.\\
\midrule
\multicolumn{2}{@{}l}{\textbf{\color{kcaccent}The Mac}}\\ \midrule
\enquote{Thonny would like to access files in your Documents folder}. & Click \menu{Allow}. (Scratch may ask the same.)\\
Scratch asks to use the microphone or camera. & Click \menu{Don't Allow}; we do not need them. You can change it later in \menu{System Settings > Privacy \& Security}.\\
Key presses do nothing in a turtle program. & Click once on the turtle window first so it has the keyboard focus, and check the program calls \texttt{screen.listen()}.\\
Screen Time blocks an app. & Add Thonny and Scratch to \enquote{Always Allowed} in \menu{System Settings > Screen Time}, or schedule sessions outside downtime.\\
\end{guidetable}

%% =====================================================================
\section{Pacing, buffers and holidays}
Each level runs \textbf{10--12 weeks}, one session a week. Weeks with no session (holidays, illness)
do not count. Choose one of two plans:
\begin{itemize}[itemsep=2pt]
  \item \textbf{12-week plan (standard):} the 12 core sessions in order.
  \item \textbf{10-week plan:} merge weeks 3 + 4 and weeks 8 + 9 as described below.
\end{itemize}
The buffers B1 (Unplugged Games Day) and B2 (Free-Build and Remix Studio) are \textbf{optional}. Use
them only when there is room inside the 12 weeks, for example a 10-week plan plus B1 as a bridge
between Scratch and Python and B2 as extra project time before the showcase.

\begin{tip}[Rules of thumb]
\begin{itemize}[itemsep=2pt]
  \item \textbf{Never skip weeks 6 and 12.} Week 6 is the review and checkpoint that makes the jump to
        Python safe; week 12 is the showcase and exit assessment that every child has been building towards.
  \item \textbf{Lose a session?} Do not add a week at the end: drop a buffer if you planned one, otherwise
        merge weeks (below), so the level still ends within 12 weeks.
  \item \textbf{One child missed a session?} Send the workbook chapter home, pair them as navigator next
        time, and use the first 5 minutes of the warm-up to catch up.
  \item \textbf{Before a holiday} end with a free-build (a natural stopping point) and no new homework.
        After a long break, start with a 10-minute review game.
\end{itemize}
\end{tip}

\subsection{How to merge weeks (10-week plan, or after lost sessions)}
\begin{itemize}[itemsep=3pt]
  \item \textbf{Weeks 3 + 4 \enquote{Catch game}:} do the week-3 warm-up and concept (events, if), then go straight
        to the week-4 build (the catch game uses events, if, a score variable and random). Week 3's own build
        becomes homework.
  \item \textbf{Weeks 8 + 9 \enquote{Loops and questions}:} teach \texttt{for}/\texttt{range} and indentation (week 8
        concept and demo), then build a polygon program that asks \enquote{How many sides?} (week 9's
        \texttt{input}/\texttt{int}). Keep week 9's 3-minute \enquote{OK to ask?} sort: the project questions in week 10 depend on it. Leave \texttt{elif} for later: spend 3 minutes on the chapter 9 flowchart at the start of week 10 and send Exercise 9.2 (Pick a Color) home as homework.
  \item Keep weeks 1, 2, 5, 6, 7, 10, 11 and 12 as they are: they carry the big transitions.
\end{itemize}

\subsection{Sample calendar for 2027 (12-week plan)}
The example uses \textbf{Saturday mornings} and the 12-week plan for all three levels. If you meet on a
weekday, shift each date to your day and skip the US holidays that fall on it (for example MLK Day,
Presidents' Day, Memorial Day, Labor Day, Indigenous Peoples' Day, Veterans Day). Spring break dates differ
by school district: move the skipped week to match yours.

\begin{center}\small\renewcommand{\arraystretch}{1.18}
\fitwidth{\begin{tabular}{@{}>{\bfseries}l l@{\hspace{1.4em}}>{\bfseries}l l@{\hspace{1.4em}}>{\bfseries}l l@{}}
\toprule
\multicolumn{2}{@{}l}{\textbf{Level 1: Jan--Apr}} & \multicolumn{2}{l}{\textbf{Level 2: May--Aug}} & \multicolumn{2}{l@{}}{\textbf{Level 3: Sep--Dec}}\\
\midrule
Jan 9 & W1 & May 8 & W1 & Sep 4 & \emph{skip: Labor Day wknd}\\
Jan 16 & W2 & May 15 & W2 & Sep 11 & W1\\
Jan 23 & W3 & May 22 & W3 & Sep 18 & W2\\
Jan 30 & W4 & May 29 & \emph{skip: Memorial Day wknd} & Sep 25 & W3\\
Feb 6 & W5 & Jun 5 & W4 & Oct 2 & W4\\
Feb 13 & \emph{skip: Presidents' Day wknd} & Jun 12 & W5 & Oct 9 & W5\\
Feb 20 & W6 checkpoint & Jun 19 & \emph{skip: Juneteenth} & Oct 16 & W6 checkpoint\\
Feb 27 & W7 & Jun 26 & W6 checkpoint & Oct 23 & W7\\
Mar 6 & W8 & Jul 3 & \emph{skip: July 4 wknd} & Oct 30 & W8\\
Mar 13 & W9 & Jul 10 & W7 & Nov 6 & W9\\
Mar 20 & W10 project kickoff & Jul 17 & W8 & Nov 13 & W10\\
Mar 27 & \emph{skip: Easter / spring break} & Jul 24 & W9 & Nov 20 & W11\\
Apr 3 & W11 & Jul 31 & W10 & Nov 27 & \emph{skip: Thanksgiving wknd}\\
Apr 10 & W12 showcase & Aug 7 & W11 & Dec 4 & W12 showcase\\
 & & Aug 14 & W12 showcase & & \\
\bottomrule
\end{tabular}}
\end{center}
\begin{itemize}[itemsep=2pt]
  \item \textbf{Level 1} has 12 sessions (Jan 9--Apr 10). MLK Day (Mon Jan 18) does not affect a Saturday
        group; if your family travels that weekend, skip Jan 16 and merge weeks 3 + 4 or 8 + 9.
  \item \textbf{Level 2} has 12 sessions (May 8--Aug 14). Want a summer break? Put it before week 9 and,
        if needed, merge weeks to keep 12 sessions or fewer.
  \item \textbf{Level 3} has 12 sessions (Sep 11--Dec 4) and ends well before winter break. Indigenous
        Peoples' Day (Mon Oct 11) and Veterans Day (Thu Nov 11) do not affect a Saturday group.
\end{itemize}

%% =====================================================================
\section{Coverage: what is taught when}

\subsection{(a) Goals across the weeks}
\cvI\ = introduced,\quad \cvR\ = revisited,\quad \cvE\ = extended. Every goal is introduced, practiced
and then extended before the exit assessment.

\begin{center}\small\renewcommand{\arraystretch}{1.35}
\fitwidth{\begin{tabular}{@{}>{\bfseries}l*{12}{c}@{}}
\toprule
Goal & W1 & W2 & W3 & W4 & W5 & W6 & W7 & W8 & W9 & W10 & W11 & W12\\
\midrule
G1 Algorithms & \cvI & \cvR & \cvR &  &  & \cvR & \cvR &  &  &  & \cvE & \\
G2 Loops &  & \cvI & \cvR & \cvR &  & \cvR &  & \cvE &  & \cvR & \cvR & \\
G3 Events, conditionals &  &  & \cvI & \cvR & \cvR & \cvR &  &  & \cvE &  & \cvR & \\
G4 Variables, random &  &  &  & \cvI & \cvR & \cvR &  &  & \cvE & \cvE & \cvR & \\
G5 Decomposition, functions &  &  &  &  & \cvI & \cvR &  &  &  & \cvE & \cvR & \cvR\\
G6 Debugging & \cvI & \cvR & \cvR & \cvR & \cvR & \cvE & \cvR & \cvR & \cvR & \cvR & \cvR & \cvR\\
G7 Pairs, explaining & \cvI & \cvR & \cvR & \cvR & \cvR & \cvR & \cvR & \cvR & \cvR & \cvR & \cvR & \cvE\\
G8 Safe, fair, accessible & \cvI & \cvR & \cvR & \cvR & \cvE & \cvR & \cvR & \cvR & \cvE & \cvR & \cvR & \cvR\\
G9 Python &  &  &  &  &  & \cvI & \cvI & \cvR & \cvR & \cvR & \cvE & \cvR\\
\bottomrule
\end{tabular}}
\end{center}
{\small Notes: G9 starts in week 6 with the Scratch-to-Python card sort and properly in week 7. G6 is
extended in week 6 (bug-hunt relay) and every week has a bug hunt. G8 is extended in week 5
(designing for accessibility) and week 9 (deciding what a program must never ask).}

\subsection{(b) Objectives and where they are taught}
The table below is generated automatically from the lesson plans (every \enquote{The student can\dots}
objective with its goal code and week).

\kcinput{generated/coverage}

\subsection{(c) The spiral map: how big ideas come back}
Level 1 lays the foundation; Levels 2 and 3 revisit and extend each idea.

\begin{guidetable}{@{}>{\bfseries\RaggedRight}p{2.6cm}>{\RaggedRight}p{\dimexpr(\linewidth-2.6cm-6\tabcolsep)/3}>{\RaggedRight}p{\dimexpr(\linewidth-2.6cm-6\tabcolsep)/3}>{\RaggedRight\arraybackslash}p{\dimexpr(\linewidth-2.6cm-6\tabcolsep)/3}@{}}
\toprule
Big idea & Level 1 & Level 2 & Level 3\\
\midrule\endhead
\bottomrule\endfoot
Sequence and algorithms & W1 \cvI\ (unplugged), W7 \cvR\ (text) & W2 \cvR\ (functions), W10 \cvE\ (pseudocode) & W4--5 \cvE\ (named algorithms, efficiency)\\
Loops & W2 \cvI\ (Scratch), W8 \cvR\ (for) & W4 \cvE\ (while, nested), W9 \cvE\ (game loop) & W4 \cvR\ (search loops), W7 \cvE\ (recursion)\\
Events and conditionals & W3 \cvI, W9 \cvR\ (if/elif) & W1 \cvE\ (and/or), W9 \cvE\ (onkey/ontimer) & W2--3 \cvR\ (methods react), project \cvE\\
Variables and data & W4 \cvI, W9 \cvR\ (input) & W1 \cvE\ (types), W3 \cvE\ (lists), W5 \cvE\ (dicts), W7 \cvE\ (files) & W2 \cvE\ (attributes), W5 \cvE\ (stacks, queues), W7 \cvE\ (grids), W8 \cvE\ (CSV)\\
Abstraction and functions & W5 \cvI\ (My Blocks), W10 \cvR\ (def) & W2 \cvE\ (return), W8 \cvR\ (testable functions) & W2--3 \cvE\ (classes), W10 \cvE\ (modules)\\
Debugging and testing & Every week (bug hunt), W6 relay & W6 \cvE\ (strategies), W8 \cvE\ (pytest) & W6 \cvE\ (debugger), project \cvE\ (5+ tests)\\
Collaboration & Pair roles every week & W10 \cvE\ (task split) & W1 \cvE\ (git), W10 \cvE\ (code review)\\
Art and games & W2 pen art, W4 catch game, W10 random art & W3 palettes, W4 pixel art, W9 games & W3 shape families, W7 fractals, project\\
Citizenship and AI ethics & Privacy, credit, kindness, accessibility & Data privacy, encryption, fair games, attribution & Git secrets, consent, bias, AI ethics, licenses\\
\end{guidetable}

\end{document}
